% Gerado por scripts/migrate_markdown_to_latex.py.
% Fonte migrada: fases/01_validacao_conceitual/docs/explicacao_metodologia_validacao.md


\chapter{Explicacao narrativa da metodologia de validacao}








Este documento explica, em linguagem de apresentacao de projeto de pesquisa, como estamos ligando os conceitos cientificos ao codigo. A ideia e construir uma historia: primeiro entendemos o que cada conceito significa na literatura; entao transformamos essa definicao em uma propriedade observavel; por conta disso criamos uma implementacao pequena; e, por fim, testamos se o resultado bate com o comportamento esperado.








O ponto central e simples: antes de medir desempenho, precisamos saber se aquilo que estamos medindo e realmente o que dizemos que e. Se chamamos uma operacao de Transformer, ela precisa ter propriedades de Transformer. Se chamamos uma tecnica de quantizacao, ela precisa representar valores com menos precisao. Se chamamos pruning de ganho de hardware, ele precisa usar algum formato ou kernel que realmente explore esparsidade.



Entao, a metodologia do projeto segue esta cadeia:



\begin{Verbatim}[fontsize=\small,breaklines=true,label=text]
conceito cientifico
-> referencia confiavel
-> definicao operacional
-> implementacao pequena e transparente
-> teste deterministico
-> classificacao do nivel de validade
\end{Verbatim}



Os quatro niveis de validade usados no projeto sao:



\begin{longtable}{@{}>{\RaggedRight\arraybackslash}p{0.273\linewidth}>{\RaggedRight\arraybackslash}p{0.273\linewidth}>{\RaggedRight\arraybackslash}p{0.273\linewidth}@{}}
\toprule
Nivel & O que responde & Exemplo \\
\midrule
\endfirsthead
\toprule
Nivel & O que responde & Exemplo \\
\midrule
\endhead
Definicao & O termo esta correto de acordo com a literatura? & Transformer como arquitetura baseada em atencao. \\
Matematica & A formula local bate com um caso pequeno calculado? & \texttt{softmax(QK\textasciicircum{}T /\allowbreak{} sqrt(d\_\allowbreak{}k))V}. \\
Algoritmo & A ferramenta implementa o comportamento esperado? & PyTorch SDPA bate com nossa attention manual. \\
Hardware & Ha evidencia de ganho fisico real? & Kernel, dtype ou formato que reduza memoria/latencia. \\
\bottomrule
\end{longtable}







Por conta disso, varias conclusoes do projeto ainda sao propositalmente cautelosas. Por exemplo, uma quantizacao pode estar numericamente correta, mas isso ainda nao quer dizer que ela gerou ganho real de hardware. Do mesmo modo, uma poda pode zerar pesos corretamente, mas se a execucao continuar densa, a conclusao e conceitual/algoritmica, nao de hardware.



\section{Capitulo 1 - Rede neural e deep learning}



\subsection{Ideia inicial}







Comecamos por rede neural porque Transformer nao aparece do nada. Um Transformer e um tipo especifico de rede neural. Entao, antes de perguntar se uma NN e Transformer, precisamos perguntar se ela pelo menos tem as caracteristicas basicas de uma rede neural: parametros, entrada, transformacoes e saida.






Deep learning entra logo depois porque, em geral, estamos falando de redes com varias etapas de transformacao. Cada camada transforma a representacao anterior em outra representacao. Entao, quando falamos de Transformer, estamos falando de uma arquitetura de deep learning que processa sequencias.



\subsection{Definicao cientifica}



A referencia principal para rede neural e deep learning e:


\begin{itemize}
\item Goodfellow, Bengio e Courville, \emph{Deep Learning}, MIT Press, 2016:
.
\end{itemize}





No nosso uso, uma rede neural e um modelo parametrizado que aplica operacoes lineares e nao lineares sobre uma entrada para produzir uma saida. Isso importa porque o nosso bloco usa pesos treinaveis, como \texttt{torch.\allowbreak{}nn.\allowbreak{}Linear}, e operacoes nao lineares, como \texttt{ReLU}.



\subsection{Por que isso importa no nosso trabalho}






Se a gente nao separa 'rede neural' de 'Transformer', corremos o risco de chamar qualquer NN de Transformer. Entao, primeiro aceitamos que o bloco atual e uma rede neural. Depois, criamos testes adicionais para saber se ela pertence a familia Transformer.



\subsection{Como implementamos}




O nucleo atual esta em \texttt{validacao.\allowbreak{}transformer.\allowbreak{}Simplified\allowbreak{}Transformer\allowbreak{}Block}. Dentro dele ha camadas lineares, normalizacao e FFN:



\begin{Verbatim}[fontsize=\small,breaklines=true,label=python]
self.q_proj = torch.nn.Linear(d_model, d_model)
self.k_proj = torch.nn.Linear(d_model, d_model)
self.v_proj = torch.nn.Linear(d_model, d_model)
self.out_proj = torch.nn.Linear(d_model, d_model)
self.norm1 = torch.nn.LayerNorm(d_model)
self.ffn = torch.nn.Sequential(
    torch.nn.Linear(d_model, ffn_hidden_dim),
    torch.nn.ReLU(),
    torch.nn.Linear(ffn_hidden_dim, d_model),
)
self.norm2 = torch.nn.LayerNorm(d_model)
\end{Verbatim}



\subsection{Como testamos}



Testamos se a rede executa com entrada sequencial e preserva shape:



\begin{Verbatim}[fontsize=\small,breaklines=true,label=python]
block = Simplified\allowbreak{}Transformer\allowbreak{}Block(d_model=4, num_heads=2, ffn_hidden_dim=8).eval()
sequence = torch.randn(2, 3, 4)
output = block(sequence)

assert output.shape == sequence.shape
\end{Verbatim}



\subsection{Com o que comparamos}





Aqui a comparacao ainda nao e com um paper de resultado. A comparacao e com a definicao operacional de uma rede neural: existe entrada, existem parametros, existe composicao de camadas e existe saida.



\subsection{O que podemos concluir}




Podemos concluir que o objeto atual e uma rede neural profunda pequena e controlada.



\subsection{O que ainda nao podemos concluir}




Ainda nao podemos concluir que ela e Transformer so por ser uma NN. Por conta disso, vem o proximo passo: validar os criterios especificos de Transformer.



\section{Capitulo 2 - Transformer}



\subsection{Ideia inicial}





Um Transformer processa sequencias usando atencao. A ideia intuitiva e: cada posicao da sequencia pode consultar outras posicoes e decidir quais informacoes devem influenciar sua representacao.






Entao, se o nosso estudo fala em eficiencia de operacoes de Transformers, a NN usada precisa ter pelo menos as pecas centrais dessa familia: entrada sequencial, Q/K/V, attention, dependencia entre posicoes, informacao posicional, residual, normalizacao e feed-forward.



\subsection{Definicao cientifica}



A referencia principal e:


\begin{itemize}
\item Vaswani et al., \emph{Attention Is All You Need}, 2017:
.
\end{itemize}



Esse trabalho define a arquitetura Transformer como uma arquitetura baseada em mecanismos de atencao. A formula central da atencao escalonada usada no paper e:



\begin{Verbatim}[fontsize=\small,breaklines=true,label=text]
Attention(Q, K, V) = softmax(QK^T / sqrt(d_k))V
\end{Verbatim}



\subsection{Por que isso importa no nosso trabalho}





Isso importa porque o nome 'Transformer' carrega uma afirmacao cientifica. Se usarmos uma NN comum e chamarmos de Transformer, o estudo perde validade. Entao, em vez de confiar no nome da classe, criamos uma auditoria arquitetural.



\subsection{Como implementamos}



A classe atual e um bloco pequeno:



\begin{Verbatim}[fontsize=\small,breaklines=true,label=python]
class Simplified\allowbreak{}Transformer\allowbreak{}Block(torch.nn.Module):
    """Bloco pequeno para validacao: MHA + residual/norm + FFN."""
\end{Verbatim}





Ela nao se declara como Transformer completo. Ela se declara como bloco simplificado. Isso e importante porque a implementacao nao tem uma pilha completa de camadas, cabeca de tarefa e treinamento final.



\subsection{Como testamos}



Usamos \texttt{validacao.\allowbreak{}auditoria\_\allowbreak{}transformer.\allowbreak{}audit\_\allowbreak{}transformer\_\allowbreak{}module}:



\begin{Verbatim}[fontsize=\small,breaklines=true,label=python]
report = audit_transformer_module(block, d_model=4, seq_len=3)

assert report.architecture_label == "bloco_transformer_simplificado"
assert report.is_transformer_family
assert not report.can_be_called_transformer
assert report.can_be_called_simplified_block
\end{Verbatim}



\subsection{Com o que comparamos}



Comparamos com uma lista de criterios derivados do Transformer original:


\begin{itemize}
\item entrada sequencial;
\item projecoes Q/K/V;
\item atencao escalonada;
\item dependencia entre posicoes;
\item informacao posicional;
\item residual, normalizacao e FFN;
\item saida deterministica em inferencia.
\end{itemize}


Tambem comparamos contra um controle negativo:



\begin{Verbatim}[fontsize=\small,breaklines=true,label=python]
plain_nn = torch.nn.Sequential(
    torch.nn.Linear(4, 4),
    torch.nn.ReLU(),
    torch.nn.Linear(4, 4),
).eval()

report = audit_transformer_module(plain_nn, d_model=4, seq_len=3)

assert report.architecture_label == "nao_aderente"
\end{Verbatim}



\subsection{O que podemos concluir}




Podemos concluir que a NN atual pertence a familia Transformer e deve ser chamada de:



\begin{Verbatim}[fontsize=\small,breaklines=true,label=text]
bloco Transformer simplificado
\end{Verbatim}



\subsection{O que ainda nao podemos concluir}





Nao podemos chamar a implementacao atual de Transformer completo. Entao, no texto do artigo, a formulacao correta e: 'avaliamos um bloco Transformer simplificado controlado'.



\subsection{Ponte direta entre Vaswani e os nossos scripts}







Entao, para nao ficar uma validacao apenas verbal, criamos uma ponte de rastreabilidade. A logica e simples: primeiro olhamos para o artigo de Vaswani, depois extraimos uma propriedade que pode aparecer no codigo, depois apontamos o script Python que implementa essa propriedade, e por fim apontamos o teste que confere se o comportamento bate com o esperado.




Essa ponte ficou documentada de forma completa em \texttt{fases/\allowbreak{}01\_\allowbreak{}validacao\_\allowbreak{}conceitual/\allowbreak{}docs/\allowbreak{}validacao\_\allowbreak{}transformer.\allowbreak{}md}. A ideia central dela e:



\begin{Verbatim}[fontsize=\small,breaklines=true,label=text]
Vaswani et al. (2017)
-> propriedade esperada
-> script Python
-> teste deterministico
-> limite da conclusao
\end{Verbatim}



Por exemplo, Vaswani define a attention escalonada como:



\begin{Verbatim}[fontsize=\small,breaklines=true,label=text]
Attention(Q, K, V) = softmax(QK^T / sqrt(d_k))V
\end{Verbatim}




Entao, por conta disso, o nosso script \texttt{fases/\allowbreak{}01\_\allowbreak{}validacao\_\allowbreak{}conceitual/\allowbreak{}validacao/\allowbreak{}attention.\allowbreak{}py} precisa fazer essa conta explicitamente:



\begin{Verbatim}[fontsize=\small,breaklines=true,label=python]
scores = torch.matmul(query, key.transpose(-1, -2)) / math.sqrt(d_k)
weights = torch.softmax(scores, dim=-1)
return torch.matmul(weights, value)
\end{Verbatim}







Mas ainda nao basta o codigo 'parecer' correto. Por conta disso, o teste \texttt{test\_\allowbreak{}scaled\_\allowbreak{}dot\_\allowbreak{}product\_\allowbreak{}attention\_\allowbreak{}matches\_\allowbreak{}independent\_\allowbreak{}numpy\_\allowbreak{}reference} compara a saida com um caso pequeno calculado de forma independente, e o teste \texttt{test\_\allowbreak{}scaled\_\allowbreak{}dot\_\allowbreak{}product\_\allowbreak{}attention\_\allowbreak{}matches\_\allowbreak{}pytorch\_\allowbreak{}reference} compara com \texttt{torch.\allowbreak{}nn.\allowbreak{}functional.\allowbreak{}scaled\_\allowbreak{}dot\_\allowbreak{}product\_\allowbreak{}attention}.



Entao, para esse trecho especifico, podemos dizer:



\begin{Verbatim}[fontsize=\small,breaklines=true,label=text]
A nossa scaled dot-product attention bate com a formula de Vaswani e com uma
implementacao de referencia do PyTorch.
\end{Verbatim}





Agora, isso ainda nao autoriza dizer que temos uma Transformer completa. Isso autoriza dizer que temos o trecho de attention validado. Por isso fazemos a mesma coisa para as outras pecas:



\begin{longtable}{@{}>{\RaggedRight\arraybackslash}p{0.273\linewidth}>{\RaggedRight\arraybackslash}p{0.273\linewidth}>{\RaggedRight\arraybackslash}p{0.273\linewidth}@{}}
\toprule
Trecho do artigo & O que exigimos no codigo & Como testamos \\
\midrule
\endfirsthead
\toprule
Trecho do artigo & O que exigimos no codigo & Como testamos \\
\midrule
\endhead
Entrada sequencial & Tensor \texttt{[batch, seq\_\allowbreak{}len, d\_\allowbreak{}model]}. & O bloco recebe a sequencia e preserva o shape. \\
Projecoes lineares & Operacao \texttt{Y = XW\textasciicircum{}T + b}. & Comparacao com valor manual conhecido. \\
Q, K e V & Tres projecoes separadas da sequencia. & Auditoria verifica \texttt{q\_\allowbreak{}proj}, \texttt{k\_\allowbreak{}proj} e \texttt{v\_\allowbreak{}proj}. \\
Scaled attention & \texttt{softmax(QK\textasciicircum{}T /\allowbreak{} sqrt(d\_\allowbreak{}k))V}. & Comparacao com NumPy, valor manual e PyTorch. \\
Multi-head & Separar heads, aplicar attention e recombinar. & Comparacao com SDPA aplicada por head. \\
Dependencia sequencial & Mudar um token pode mudar outra posicao. & Teste altera uma posicao e observa outra. \\
Posicao & Algum sinal posicional antes da attention. & Auditoria verifica positional encoding. \\
Residual/norm/FFN & Componentes de bloco Transformer. & Auditoria estrutural do bloco. \\
Controle negativo & NN comum deve ser rejeitada. & \texttt{Linear + ReLU + Linear} recebe \texttt{nao\_\allowbreak{}aderente}. \\
\bottomrule
\end{longtable}




Essa tabela e importante para a apresentacao porque ela mostra que nao estamos dizendo 'e Transformer porque eu chamei de Transformer'. Estamos dizendo:



\begin{Verbatim}[fontsize=\small,breaklines=true,label=text]
Este trecho pertence a familia Transformer porque a propriedade X aparece no
artigo, aparece no codigo, e passa no teste Y.
\end{Verbatim}







Com isso, a conclusao fica calibrada. Podemos afirmar que os scripts atuais validam trechos centrais de Transformer e um bloco Transformer simplificado. Ainda nao podemos afirmar que reproduzimos uma Transformer completa, porque a arquitetura completa exigiria pilha de blocos, encoder/decoder ou pelo menos um encoder completo, embeddings, cabeca de tarefa e validacao em tarefa real.



\section{Capitulo 3 - Bloco Transformer simplificado}



\subsection{Ideia inicial}





Depois de entender Transformer, a pergunta vira: qual parte do Transformer vamos estudar? Como o projeto ainda esta em fase controlada, usamos um bloco. Isso permite entender causa e efeito sem misturar muitas variaveis.




Entao, em vez de medir um modelo enorme logo de inicio, estudamos um bloco pequeno onde conseguimos auditar cada parte.



\subsection{Definicao cientifica}





O bloco se inspira no Transformer de Vaswani et al. Ele tem attention, projecoes, residual, normalizacao e FFN. Como e uma simplificacao, ele nao e uma arquitetura completa.



\subsection{Por que isso importa no nosso trabalho}






Isso importa porque a validade cientifica vem do controle. Se usarmos um modelo grande sem entender o caminho interno, talvez vejamos uma mudanca de latencia sem saber se veio da attention, da quantizacao, do kernel, da memoria ou de alguma otimizacao escondida.



\subsection{Como implementamos}



O forward do bloco mostra a historia inteira:



\begin{Verbatim}[fontsize=\small,breaklines=true,label=python]
hidden = self.add_positional_information(sequence)
query, key, value = self.project_qkv(hidden)
attention = self.attention_core(query, key, value)
attention = combine_heads(attention)

hidden = self.norm1(hidden + self.out_proj(attention))
return self.norm2(hidden + self.ffn(hidden))
\end{Verbatim}



Entao, o fluxo e:



\begin{Verbatim}[fontsize=\small,breaklines=true,label=text]
sequencia
-> informacao posicional
-> Q/K/V
-> attention
-> recombinacao das heads
-> residual + norm
-> FFN
-> residual + norm
\end{Verbatim}



\subsection{Como testamos}



O teste confere componentes e shape:



\begin{Verbatim}[fontsize=\small,breaklines=true,label=python]
assert output.shape == sequence.shape
assert hasattr(block, "q_proj")
assert hasattr(block, "k_proj")
assert hasattr(block, "v_proj")
assert hasattr(block, "norm1")
assert hasattr(block, "ffn")
\end{Verbatim}



E a auditoria confere o comportamento:



\begin{Verbatim}[fontsize=\small,breaklines=true,label=python]
assert report.has_qkv_projection
assert report.has_scaled_attention
assert report.has_cross_position_dependency
assert report.has_positional_information_or_justification
assert report.has_residual_norm_ffn_or_justification
\end{Verbatim}



\subsection{Com o que comparamos}




Comparamos com os criterios arquiteturais do Transformer e com a API de referencia do PyTorch para a attention interna.



\subsection{O que podemos concluir}



Podemos concluir que o bloco e adequado para o recorte experimental.



\subsection{O que ainda nao podemos concluir}




Ainda nao podemos concluir desempenho final de um Transformer real em producao, porque isso exigiria arquitetura completa, tarefa, dataset e benchmark maior.



\section{Capitulo 4 - Self-attention}



\subsection{Ideia inicial}




Self-attention e a ideia de uma sequencia olhar para ela mesma. Entao, cada posicao calcula o quanto deve considerar outras posicoes.




Isso importa porque, sem dependencia entre posicoes, uma operacao nao esta capturando o aspecto sequencial que torna Transformers interessantes.



\subsection{Definicao cientifica}




A referencia e Vaswani et al. (2017). No nosso recorte, self-attention significa que Q, K e V sao derivados da mesma sequencia.



\subsection{Por que isso importa no nosso trabalho}





Se mudamos um token e nenhuma outra posicao muda, entao a operacao nao esta misturando informacao da sequencia. Por conta disso, testamos explicitamente se alterar uma posicao pode alterar outra.



\subsection{Como implementamos}



A attention central usa Q, K e V:



\begin{Verbatim}[fontsize=\small,breaklines=true,label=python]
scores = torch.matmul(query, key.transpose(-1, -2)) / math.sqrt(d_k)
weights = torch.softmax(scores, dim=-1)
return torch.matmul(weights, value)
\end{Verbatim}



\subsection{Como testamos}



O teste cria duas versoes de \texttt{value}, alterando a segunda posicao:



\begin{Verbatim}[fontsize=\small,breaklines=true,label=python]
query = torch.zeros(1, 2, 2)
key = torch.zeros(1, 2, 2)
value_a = torch.tensor([[[1.0, 1.0], [3.0, 3.0]]])
value_b = torch.tensor([[[1.0, 1.0], [9.0, 9.0]]])

output_a = scaled_dot_product_attention_torch(query, key, value_a)
output_b = scaled_dot_product_attention_torch(query, key, value_b)

assert not torch.allclose(output_a[:, 0, :], output_b[:, 0, :])
\end{Verbatim}



\subsection{Com o que comparamos}




Comparamos com a propriedade esperada de self-attention: uma posicao pode depender de outra posicao da mesma sequencia.



\subsection{O que podemos concluir}



Podemos concluir que a operacao permite dependencia entre posicoes.



\subsection{O que ainda nao podemos concluir}




Isso nao prova qualidade de modelagem, nem desempenho, nem acuracia em tarefa. Prova apenas a propriedade conceitual necessaria.



\section{Capitulo 5 - Scaled dot-product attention}



\subsection{Ideia inicial}





Scaled dot-product attention e a conta central. Ela calcula similaridade entre queries e keys, normaliza essas similaridades e usa o resultado para combinar values.




Entao, esse e um dos testes mais importantes: se essa formula estiver errada, o resto do estudo perde base.



\subsection{Definicao cientifica}



A formula vem de Vaswani et al.:



\begin{Verbatim}[fontsize=\small,breaklines=true,label=text]
Attention(Q, K, V) = softmax(QK^T / sqrt(d_k))V
\end{Verbatim}



\subsection{Por que isso importa no nosso trabalho}





Essa formula e o nucleo matematico da attention usada no bloco. Por conta disso, nao basta confiar no PyTorch. Implementamos uma versao manual e comparamos com uma referencia independente.



\subsection{Como implementamos}



Em PyTorch:



\begin{Verbatim}[fontsize=\small,breaklines=true,label=python]
d_k = query.shape[-1]
scores = torch.matmul(query, key.transpose(-1, -2)) / math.sqrt(d_k)
weights = torch.softmax(scores, dim=-1)
return torch.matmul(weights, value)
\end{Verbatim}



Em NumPy, a mesma ideia:



\begin{Verbatim}[fontsize=\small,breaklines=true,label=python]
scores = np.matmul(query, np.swapaxes(key, -1, -2)) / math.sqrt(d_k)
weights = softmax_numpy(scores, axis=-1)
return np.matmul(weights, value)
\end{Verbatim}



\subsection{Como testamos}



Usamos tensores pequenos:



\begin{Verbatim}[fontsize=\small,breaklines=true,label=python]
query = torch.tensor([[[1.0, 0.0], [0.0, 1.0]]])
key = torch.tensor([[[1.0, 0.0], [0.0, 1.0]]])
value = torch.tensor([[[1.0, 2.0], [3.0, 4.0]]])
\end{Verbatim}



E exigimos o resultado esperado pela formula:



\begin{Verbatim}[fontsize=\small,breaklines=true,label=python]
expected_from_formula = torch.tensor(
    [[[1.6604769, 2.6604769], [2.3395231, 3.3395231]]]
)

assert torch.allclose(actual, expected_from_formula, atol=1e-6)
\end{Verbatim}



\subsection{Com o que comparamos}



Comparamos com:


\begin{enumerate}
\item calculo NumPy independente;
\item resultado numerico derivado da formula;
\item \texttt{torch.\allowbreak{}nn.\allowbreak{}functional.\allowbreak{}scaled\_\allowbreak{}dot\_\allowbreak{}product\_\allowbreak{}attention};
\item \texttt{keras.\allowbreak{}ops.\allowbreak{}nn.\allowbreak{}dot\_\allowbreak{}product\_\allowbreak{}attention} com backend TensorFlow, usando os mesmos
\texttt{Q}, \texttt{K}, \texttt{V} e mascara.
\end{enumerate}


\subsection{O que podemos concluir}






Podemos concluir que a attention manual esta matematicamente aderente ao conceito e numericamente compativel com APIs publicas de dois frameworks. A execucao canônica aprovou 9/9 comparacoes, com erro absoluto maximo de \texttt{2.\allowbreak{}38418579e-07} para \texttt{atol=1e-6} e \texttt{rtol=1e-5}.



\subsection{O que ainda nao podemos concluir}




Attention correta nao significa Transformer completo. Ela e uma parte necessaria, mas nao suficiente.



\section{Capitulo 6 - Multi-head attention}



\subsection{Ideia inicial}





Multi-head attention divide a representacao em varias heads. Entao, em vez de uma unica attention olhar para tudo de uma forma, varias attentions menores processam partes da representacao.



\subsection{Definicao cientifica}




A referencia continua sendo Vaswani et al. (2017). A arquitetura Transformer usa varias heads para projetar, aplicar attention e recombinar informacoes.



\subsection{Por que isso importa no nosso trabalho}





Se o bloco usa multi-head attention, precisamos saber se o split e o combine estao corretos. Por conta disso, validamos shapes e comparamos com a referencia do PyTorch.



\subsection{Como implementamos}



O split:



\begin{Verbatim}[fontsize=\small,breaklines=true,label=python]
return tensor.reshape(batch_size, seq_len, num_heads, head_dim).transpose(1, 2)
\end{Verbatim}



O combine:



\begin{Verbatim}[fontsize=\small,breaklines=true,label=python]
return tensor.transpose(1, 2).contiguous().reshape(
    batch_size,
    seq_len,
    num_heads * head_dim,
)
\end{Verbatim}



E a attention multi-head transparente:



\begin{Verbatim}[fontsize=\small,breaklines=true,label=python]
query_heads = split_heads(query, num_heads)
key_heads = split_heads(key, num_heads)
value_heads = split_heads(value, num_heads)
attention_heads = scaled_dot_product_attention_torch(query_heads, key_heads, value_heads)
return combine_heads(attention_heads)
\end{Verbatim}



\subsection{Como testamos}



Testamos shape e equivalencia:



\begin{Verbatim}[fontsize=\small,breaklines=true,label=python]
sequence = torch.randn(2, 3, 4)
num_heads = 2

actual = multi_head_attention_from_qkv(
    sequence,
    sequence,
    sequence,
    num_heads=num_heads,
)

assert actual.shape == sequence.shape
\end{Verbatim}



\subsection{Com o que comparamos}




Comparamos com \texttt{torch.\allowbreak{}nn.\allowbreak{}functional.\allowbreak{}scaled\_\allowbreak{}dot\_\allowbreak{}product\_\allowbreak{}attention} aplicada nas heads ja separadas.



\subsection{O que podemos concluir}



Podemos concluir que a mecanica multi-head esta correta no caso pequeno.



\subsection{O que ainda nao podemos concluir}




Nao podemos concluir que o numero de heads escolhido e otimo. Isso seria uma pergunta experimental diferente.



\section{Capitulo 7 - Projecao densa}



\subsection{Ideia inicial}





Projecao densa e uma camada linear. Ela pega um vetor e transforma em outro por meio de uma matriz de pesos. Entao, antes de Q, K e V existirem, precisamos projetar a entrada.



\subsection{Definicao cientifica}



A operacao e:



\begin{Verbatim}[fontsize=\small,breaklines=true,label=text]
Y = XW^T + b
\end{Verbatim}




Ela aparece em redes neurais em Goodfellow et al. e no Transformer de Vaswani et al., onde projecoes aprendidas geram Q, K, V e saidas de attention.



\subsection{Por que isso importa no nosso trabalho}




Se a projecao densa estiver errada, Q/K/V tambem ficam errados. Entao, validamos essa operacao separadamente antes de confiar no bloco.



\subsection{Como implementamos}



\begin{Verbatim}[fontsize=\small,breaklines=true,label=python]
def dense_projection_torch(inputs, weight, bias=None):
    output = torch.matmul(inputs, weight.transpose(-1, -2))
    if bias is not None:
        output = output + bias
    return output
\end{Verbatim}



\subsection{Como testamos}



Usamos um exemplo pequeno:



\begin{Verbatim}[fontsize=\small,breaklines=true,label=python]
inputs = torch.tensor([[[1.0, 2.0], [3.0, 4.0]]])
weight = torch.tensor([[1.0, 0.0], [0.0, 2.0]])
bias = torch.tensor([0.5, -0.5])

expected = torch.tensor([[[1.5, 3.5], [3.5, 7.5]]])

assert torch.allclose(actual, expected)
\end{Verbatim}



\subsection{Com o que comparamos}



Comparamos com o resultado manual e com uma implementacao NumPy independente.



\subsection{O que podemos concluir}




Podemos concluir que a projecao densa do projeto segue a formula linear esperada.



\subsection{O que ainda nao podemos concluir}




Projecao densa sozinha nao e Transformer. Ela e uma operacao auxiliar usada dentro do bloco.



\section{Capitulo 8 - Informacao posicional}



\subsection{Ideia inicial}





Attention pura enxerga relacoes entre vetores, mas nao sabe automaticamente a ordem dos tokens. Entao, Transformers usam alguma forma de informacao posicional para indicar que a posicao 1 vem antes da posicao 2.



\subsection{Definicao cientifica}





Vaswani et al. usam positional encoding no Transformer original. No nosso recorte, usamos encoding sinusoidal, que adiciona sinais dependentes da posicao aos vetores da sequencia.



\subsection{Por que isso importa no nosso trabalho}





Se removemos informacao posicional sem justificativa, o bloco fica incompleto para ser chamado de bloco Transformer. Por conta disso, a auditoria verifica explicitamente se existe informacao posicional ou justificativa formal.



\subsection{Como implementamos}



\begin{Verbatim}[fontsize=\small,breaklines=true,label=python]
position = torch.arange(seq_len, device=device, dtype=dtype).unsqueeze(1)
div_term = torch.exp(
    torch.arange(0, d_model, 2, device=device, dtype=dtype)
    * (-math.log(10000.0) / d_model)
)
encoding[:, 0::2] = torch.sin(position * div_term)
encoding[:, 1::2] = torch.cos(position * div_term[: encoding[:, 1::2].shape[1]])
\end{Verbatim}



E no bloco:



\begin{Verbatim}[fontsize=\small,breaklines=true,label=python]
hidden = self.add_positional_information(sequence)
\end{Verbatim}



\subsection{Como testamos}



A auditoria verifica:



\begin{Verbatim}[fontsize=\small,breaklines=true,label=python]
assert report.has_positional_information_or_justification
\end{Verbatim}



\subsection{Com o que comparamos}




Comparamos com o criterio arquitetural do Transformer: a sequencia precisa ter alguma informacao de posicao, ou a ausencia precisa ser justificada.



\subsection{O que podemos concluir}



Podemos concluir que o bloco atual tem informacao posicional.



\subsection{O que ainda nao podemos concluir}




Nao podemos concluir que esse encoding e o melhor para qualquer tarefa. Ele e adequado para validar o conceito.



\section{Capitulo 9 - Inferencia}



\subsection{Ideia inicial}




Inferencia e usar o modelo para produzir saidas, sem atualizar pesos. Entao, quando medimos custo computacional, queremos medir execucao, nao treinamento.



\subsection{Definicao cientifica}




Goodfellow et al. separam o uso de modelos para predicao da fase de treinamento. No PyTorch, operacionalmente usamos \texttt{eval()} e \texttt{torch.\allowbreak{}inference\_\allowbreak{}mode()}.



\subsection{Por que isso importa no nosso trabalho}




Se os pesos mudam durante o experimento, a comparacao com baseline deixa de ser justa. Entao, fixamos seed, colocamos modelo em \texttt{eval()} e usamos inferencia.



\subsection{Como implementamos}



Na auditoria:



\begin{Verbatim}[fontsize=\small,breaklines=true,label=python]
module.eval()
with torch.inference_mode():
    return module(sequence)
\end{Verbatim}



No benchmark controlado:



\begin{Verbatim}[fontsize=\small,breaklines=true,label=python]
with torch.inference_mode():
    baseline_output = baseline_model(inputs).detach()
\end{Verbatim}



\subsection{Como testamos}



A auditoria roda a mesma entrada duas vezes:



\begin{Verbatim}[fontsize=\small,breaklines=true,label=python]
first_output = _safe_forward(module, sequence)
second_output = _safe_forward(module, sequence)

deterministic_in_eval = torch.allclose(first_output, second_output, atol=atol)
\end{Verbatim}



\subsection{Com o que comparamos}



Comparamos uma saida contra a outra, para verificar determinismo.



\subsection{O que podemos concluir}




Podemos concluir que o bloco se comporta de forma deterministica em modo de avaliacao.



\subsection{O que ainda nao podemos concluir}




Isso nao garante reproducibilidade universal em qualquer hardware e qualquer kernel. Por conta disso, o benchmark registra ambiente, seed e dispositivo.



\section{Capitulo 10 - KV cache}



\subsection{Ideia inicial}





KV cache e uma economia de recomputacao. Em geracao autoregressiva, quando o modelo gera token por token, as chaves e valores dos tokens antigos nao precisam ser recalculados toda vez.




Entao, o cache guarda K e V passados. A cada novo token, calculamos K/V apenas do token novo e juntamos com o cache.



\subsection{Definicao cientifica}




O conceito aparece em inferencia de Transformers autoregressivos e em trabalhos de compressao de KV cache. Usamos como referencia moderna TurboQuant:


\begin{itemize}
\item Zandieh et al., \emph{TurboQuant}, 2025:
.
\item Google Research, *TurboQuant: Redefining AI efficiency with extreme
compression*:
.
\end{itemize}


\subsection{Por que isso importa no nosso trabalho}




KV cache e importante para memoria e latencia em inferencia. Mas, antes de medir latencia, precisamos saber se o cache preserva a mesma saida.



\subsection{Como implementamos}



O caminho com cache concatena K/V antigos com K/V do token novo:



\begin{Verbatim}[fontsize=\small,breaklines=true,label=python]
key_cache = key if key_cache is None else torch.cat([key_cache, key], dim=1)
value_cache = value if value_cache is None else torch.cat([value_cache, value], dim=1)
outputs.append(scaled_dot_product_attention_torch(query, key_cache, value_cache))
\end{Verbatim}



\subsection{Como testamos}



Comparamos atencao causal completa contra atencao incremental com cache:



\begin{Verbatim}[fontsize=\small,breaklines=true,label=python]
full_output = causal_self_attention_full(sequence, identity, identity, identity)
cached = causal_self_attention_with_cache(sequence, identity, identity, identity)

assert torch.allclose(cached.output, full_output, atol=1e-6)
assert cached.cached_key_value_tokens == 4
assert cached.naive_key_value_tokens == 10
assert cached.reused_key_value_tokens == 6
\end{Verbatim}



\subsection{Com o que comparamos}





Comparamos a saida com e sem cache. Se as saidas batem, o cache preserva o comportamento. Tambem comparamos a quantidade de K/V calculados: 4 contra 10 no caso pequeno.



\subsection{O que podemos concluir}



Podemos concluir equivalencia algoritmica e reducao de recomputacao no exemplo.



\subsection{O que ainda nao podemos concluir}




Ainda nao podemos concluir ganho de hardware. Para isso, precisamos medir memoria, latencia e throughput em execucao controlada.



\section{Capitulo 11 - Pruning e sparsity}



\subsection{Ideia inicial}





Pruning e remover ou zerar partes do modelo. A intencao e criar esparsidade: muitos zeros ou conexoes removidas. Entao, teoricamente, haveria menos trabalho a fazer.





Mas existe um detalhe importante: zerar pesos nao significa necessariamente que o hardware vai pular esses pesos. Se a multiplicacao continuar densa, a conta ainda pode custar quase a mesma coisa.



\subsection{Definicao cientifica}



Usamos:


\begin{itemize}
\item Tang et al., \emph{A Survey on Transformer Compression}, 2024:
.
\item PyTorch pruning tutorial:
.
\end{itemize}


\subsection{Por que isso importa no nosso trabalho}





Como vamos comparar pruning com quantizacao, precisamos classificar pruning sem exagero. Entao, se o pruning apenas cria mascara, classificamos como conceitual/algoritmico.



\subsection{Como implementamos}



A observacao de sparsity conta zeros:



\begin{Verbatim}[fontsize=\small,breaklines=true,label=python]
total = int(tensor.numel())
zeros = int((tensor == 0).sum().item())
ratio = 0.0 if total == 0 else zeros / total
\end{Verbatim}




E a classificacao considera hardware apenas se houver armazenamento e kernel esparsos:



\begin{Verbatim}[fontsize=\small,breaklines=true,label=python]
if self.uses_sparse_storage and self.uses_sparse_kernel:
    return "hardware"
return "conceitual_algoritmico"
\end{Verbatim}



\subsection{Como testamos}



Usamos uma camada com 8 pesos e aplicamos pruning de 50\%:



\begin{Verbatim}[fontsize=\small,breaklines=true,label=python]
prune.l1_unstructured(layer, name="weight", amount=0.5)

observation = observe_sparsity(
    layer.weight,
    uses_sparse_storage=False,
    uses_sparse_kernel=False,
)

assert observation.zero_elements == 4
assert observation.sparsity_ratio == pytest.approx(0.5)
assert observation.validity_level == "conceitual_algoritmico"
assert hasattr(layer, "weight_mask")
\end{Verbatim}



\subsection{Com o que comparamos}




Comparamos com a expectativa operacional: 50\% de 8 pesos = 4 zeros. Tambem verificamos a presenca de \texttt{weight\_\allowbreak{}mask}, coerente com o mecanismo do PyTorch.



\subsection{O que podemos concluir}



Podemos concluir que o pruning cria esparsidade conceitual/algoritmica.



\subsection{O que ainda nao podemos concluir}




Nao podemos dizer que houve ganho real de hardware, porque o teste nao usa storage esparso nem kernel esparso.



\section{Capitulo 12 - Quantizacao}



\subsection{Ideia inicial}




Quantizacao e representar numeros com menos precisao. Entao, em vez de guardar tudo em \texttt{float32}, podemos representar pesos em \texttt{int8}, por exemplo.





Isso economiza memoria representacional, mas introduz erro numerico. Por conta disso, precisamos medir duas coisas: se a representacao realmente ficou menor e se o erro ficou controlado.



\subsection{Definicao cientifica}



Usamos:


\begin{itemize}
\item Tang et al. (2024), sobre compressao de Transformers.
\item Zandieh et al. (2025), TurboQuant.
\item Documentacao \texttt{torchao}:
.
\end{itemize}


\subsection{Por que isso importa no nosso trabalho}





Quantizacao e uma das tecnicas centrais do estudo. Mas se quantizamos e logo voltamos para \texttt{float32}, isso e apenas simulacao numerica. Entao, separamos validade numerica de validade de hardware.



\subsection{Como implementamos}



A quantizacao manual simetrica:



\begin{Verbatim}[fontsize=\small,breaklines=true,label=python]
max_abs = tensor.abs().max()
scale = max_abs / 127.0
quantized = torch.clamp(torch.round(tensor / scale), -127, 127).to(torch.int8)
dequantized = quantized.to(torch.float32) * scale
\end{Verbatim}



E com \texttt{torchao} v2:



\begin{Verbatim}[fontsize=\small,breaklines=true,label=python]
from torchao.quantization import Int8WeightOnlyConfig, quantize_

quantize_(model, Int8WeightOnlyConfig(version=2))
\end{Verbatim}



\subsection{Como testamos}



Para a quantizacao manual:



\begin{Verbatim}[fontsize=\small,breaklines=true,label=python]
tensor = torch.tensor([-2.0, -1.0, 0.0, 1.0, 2.0], dtype=torch.float32)
quantized, scale, dequantized = symmetric_int8_quantize(tensor)

assert quantized.dtype == torch.int8
assert q_observation.storage_bytes == 5
assert tensor.numel() * tensor.element_size() == 20
assert torch.max(torch.abs(tensor - dequantized)).item() <= scale.item()
\end{Verbatim}



Para \texttt{torchao}:



\begin{Verbatim}[fontsize=\small,breaklines=true,label=python]
with warnings.catch_warnings(record=True) as caught:
    warnings.simplefilter("always")
    quantize_(model, Int8WeightOnlyConfig(version=2))

assert caught == []
assert type(model.weight).__name__ == "Int8Tensor"
assert model.weight.qdata.dtype == torch.int8
\end{Verbatim}



\subsection{Com o que comparamos}



Comparamos:


\begin{itemize}
\item \texttt{float32}: 5 valores * 4 bytes = 20 bytes;
\item \texttt{int8}: 5 valores * 1 byte = 5 bytes;
\item erro maximo contra o limite dado pela escala;
\item representacao interna \texttt{qdata} em \texttt{int8} no \texttt{torchao}.
\end{itemize}


\subsection{O que podemos concluir}



Podemos concluir validade numerica e algoritmica da quantizacao no recorte.



\subsection{O que ainda nao podemos concluir}




Ainda nao podemos dizer 'ganho de hardware' sem demonstrar kernel/caminho de execucao compativel.



\section{Capitulo 13 - Latencia, throughput, memoria e FLOPs}



\subsection{Ideia inicial}






Depois que os conceitos estao validados, queremos medir custo. Mas custo nao e uma coisa so. Latencia e tempo por execucao. Throughput e quantidade processada por tempo. Memoria e armazenamento usado. FLOPs e uma estimativa teorica de operacoes.




Entao, se misturarmos tudo, a conclusao fica confusa. Por conta disso, medimos cada dimensao separadamente.



\subsection{Definicao cientifica}





Usamos Hennessy e Patterson, \emph{Computer Architecture: A Quantitative Approach}, como referencia geral para avaliacao quantitativa de desempenho, memoria, latencia e throughput.




Tambem usamos Vaswani et al. para entender o custo de operacoes de attention em sequencias.



\subsection{Por que isso importa no nosso trabalho}





Uma tecnica pode reduzir memoria e nao reduzir latencia. Outra pode reduzir FLOPs teoricos e nao acelerar no hardware real. Entao, cada resultado precisa dizer exatamente qual metrica mudou.



\subsection{Como implementamos FLOPs}



Para projecao densa:



\begin{Verbatim}[fontsize=\small,breaklines=true,label=python]
matmul_flops = 2 * batch_size * seq_len * in_features * out_features
bias_flops = batch_size * seq_len * out_features if bias else 0
return int(matmul_flops + bias_flops)
\end{Verbatim}



Para attention:



\begin{Verbatim}[fontsize=\small,breaklines=true,label=python]
score_flops = 2 * batch_size * num_heads * seq_len * seq_len * head_dim
weighted_value_flops = 2 * batch_size * num_heads * seq_len * seq_len * head_dim
return int(score_flops + weighted_value_flops)
\end{Verbatim}



\subsection{Como testamos FLOPs}



\begin{Verbatim}[fontsize=\small,breaklines=true,label=python]
assert dense_projection_flops(
    batch_size=2,
    seq_len=3,
    in_features=4,
    out_features=5,
) == 270

assert attention_matmul_flops(
    batch_size=2,
    seq_len=3,
    d_model=4,
    num_heads=2,
) == 288
\end{Verbatim}



\subsection{Como implementamos benchmark controlado}



O runner registra ambiente, seed, warmup, repeticoes e sincronizacao:



\begin{Verbatim}[fontsize=\small,breaklines=true,label=python]
metadata = {
    **env,
    "warmup": config.warmup,
    "repetitions": config.repetitions,
    "scenarios": list(config.scenarios),
    "used_cuda_synchronization": used_cuda_synchronization,
}
\end{Verbatim}



\subsection{Com o que comparamos}




FLOPs sao comparados com formulas teoricas. Latencia, throughput e memoria serao comparados contra baseline no CSV gerado pelo benchmark.



\subsection{O que podemos concluir}



Podemos concluir que temos um protocolo de medicao pronto e validado por teste.



\subsection{O que ainda nao podemos concluir}




Ainda nao podemos concluir desempenho final antes de rodar benchmarks amplos e controlados.



\section{Capitulo 14 - MSE, MAE, R2 e similaridade de cosseno}



\subsection{Ideia inicial}




Quando aplicamos pruning ou quantizacao, a saida pode mudar. Entao, precisamos medir quanto ela mudou em relacao ao baseline.




Por conta disso usamos quatro metricas: MSE, MAE, R2 e similaridade de cosseno. Elas olham para erro e preservacao da direcao da representacao.



\subsection{Definicao cientifica}



Usamos:


\begin{itemize}
\item Hastie, Tibshirani e Friedman, \emph{The Elements of Statistical Learning}:
.
\item scikit-learn MSE:
.
\item scikit-learn MAE:
.
\item scikit-learn R2:
.
\item scikit-learn cosine similarity:
.
\end{itemize}


\subsection{Por que isso importa no nosso trabalho}




Se uma tecnica acelera mas destroi a saida, ela pode nao ser util. Entao, desempenho precisa ser lido junto com fidelidade numerica.



\subsection{Como implementamos}



MSE:



\begin{Verbatim}[fontsize=\small,breaklines=true,label=python]
return float(np.mean(np.square(reference_arr - candidate_arr)))
\end{Verbatim}



MAE:



\begin{Verbatim}[fontsize=\small,breaklines=true,label=python]
return float(np.mean(np.abs(reference_arr - candidate_arr)))
\end{Verbatim}



R2:



\begin{Verbatim}[fontsize=\small,breaklines=true,label=python]
return float(1.0 - (ss_res / ss_tot))
\end{Verbatim}



Cosseno:



\begin{Verbatim}[fontsize=\small,breaklines=true,label=python]
return float(np.dot(reference_arr, candidate_arr) / (reference_norm * candidate_norm))
\end{Verbatim}



\subsection{Como testamos}



Para MSE, MAE e R2:



\begin{Verbatim}[fontsize=\small,breaklines=true,label=python]
reference = np.array([1.0, 2.0, 3.0])
candidate = np.array([1.0, 2.0, 4.0])

assert mean_squared_error(reference, candidate) == pytest.approx(1.0 / 3.0)
assert mean_absolute_error(reference, candidate) == pytest.approx(1.0 / 3.0)
assert r2_score(reference, candidate) == pytest.approx(0.5)
\end{Verbatim}



Para cosseno:



\begin{Verbatim}[fontsize=\small,breaklines=true,label=python]
reference = np.array([1.0, 1.0, 0.0])
candidate = np.array([1.0, 0.0, 0.0])

expected = 1.0 / math.sqrt(2.0)
assert cosine_similarity(reference, candidate) == pytest.approx(expected)
\end{Verbatim}



\subsection{Com o que comparamos}



Comparamos com valores manuais conhecidos e com \texttt{scikit-learn}.



\subsection{O que podemos concluir}




Podemos concluir que as metricas estao corretas para comparar saidas contra o baseline.



\subsection{O que ainda nao podemos concluir}




As metricas nao dizem, sozinhas, se uma tecnica e boa. Elas precisam ser interpretadas junto com latencia, memoria, throughput e nivel de validade.



\section{Capitulo 15 - Protocolo experimental e matriz de validade}



\subsection{Ideia inicial}





Agora juntamos tudo. Se cada conceito tem teste, ainda precisamos garantir que os resultados futuros sejam registrados de forma comparavel. Entao, criamos um schema CSV e uma matriz de validacao.



\subsection{Definicao cientifica}





Aqui a referencia e metodologica: experimentos precisam de rastreabilidade. No nosso caso, rastreabilidade significa ligar cada conclusao a uma fonte, a um teste e a um resultado.



\subsection{Por que isso importa no nosso trabalho}




Sem matriz, podemos acabar escrevendo conclusoes mais fortes do que os dados permitem. Entao, o pipeline obriga cada cenario a ter nivel de validade.



\subsection{Como implementamos}



O schema CSV obrigatorio:



\begin{Verbatim}[fontsize=\small,breaklines=true,label=python]
CSV_COLUMNS = [
    "torch_version",
    "cuda_version",
    "gpu_name",
    "device",
    "model_kind",
    "scenario",
    "batch_size",
    "seq_len",
    "d_model",
    "num_heads",
    "dtype",
    "pruning_sparsity",
    "quantization_method",
    "validity_level",
    "latency_ms_mean",
    "latency_ms_p50",
    "latency_ms_p95",
    "throughput_tokens_s",
    "max_memory_bytes",
    "theoretical_flops",
    "mse",
    "mae",
    "r2",
    "cosine_similarity",
]
\end{Verbatim}



E a validacao rejeita colunas faltando, colunas extras e niveis invalidos.



\subsection{Como testamos}



\begin{Verbatim}[fontsize=\small,breaklines=true,label=python]
validation = validate_benchmark_record(record)
invalid = validate_benchmark_record({"validity_level": "promessa_sem_evidencia"})

assert validation.is_valid
assert not invalid.is_valid
\end{Verbatim}



Tambem testamos que cenarios nao registrados sao rejeitados:



\begin{Verbatim}[fontsize=\small,breaklines=true,label=python]
with pytest.raises(ScenarioValidationError, match="cenario sem matriz"):
    validate_scenarios_registered(("baseline", "cenario_sem_evidencia"))
\end{Verbatim}



\subsection{Com o que comparamos}



Comparamos cada resultado contra o baseline e contra a matriz de validade.



\subsection{O que podemos concluir}




Podemos concluir que o estudo tem um caminho metodologico controlado para aceitar ou rejeitar resultados.



\subsection{O que ainda nao podemos concluir}





Ainda nao podemos concluir qual tecnica vence no desempenho final. O protocolo so garante que, quando os benchmarks forem rodados, a conclusao sera rastreavel.



\section{Fechamento - A frase metodologica do projeto}



A historia inteira pode ser apresentada assim:



\begin{icquote}
O estudo nao mede desempenho antes de validar significado. Primeiro definimos cada conceito pela literatura. Entao transformamos essa definicao em uma propriedade observavel no codigo. Depois testamos essa propriedade com casos pequenos, deterministas e comparaveis com formulas ou bibliotecas de referencia. Por conta disso, quando formos medir pruning, quantizacao, latencia, memoria e throughput, saberemos exatamente o nivel de validade de cada conclusao.
\end{icquote}








Essa frase protege o trabalho de uma fragilidade comum: usar palavras fortes sem evidencia proporcional. Entao, quando dissermos 'bloco Transformer simplificado', sabemos quais criterios ele passou. Quando dissermos 'quantizacao numerica', sabemos que houve \texttt{int8}, escala e erro medido. Quando dissermos 'pruning conceitual/algoritmico', sabemos que ha zeros e mascara, mas nao necessariamente ganho de hardware.



O estado atual dos testes confirma essa base:



\begin{Verbatim}[fontsize=\small,breaklines=true,label=powershell]
.\.venv\Scripts\python.exe -m pytest -q -W error
\end{Verbatim}



Resultado observado no ambiente de validacao cruzada em 24/08/2026:



\begin{Verbatim}[fontsize=\small,breaklines=true,label=text]
32 passed, 1 skipped
\end{Verbatim}




O teste ignorado exige CUDA. Os testes NumPy, PyTorch e TensorFlow/Keras passaram sem \texttt{skip}.






Entao, o proximo passo do projeto nao e 'inventar a metodologia'; a metodologia ja esta desenhada. O proximo passo e usar esse protocolo para rodar benchmarks controlados e interpretar os resultados sem ultrapassar o nivel de evidencia observado.
